%======================================================================
\section{Heuristic: from ranks to returns}
\label{sec:heuristic}

The construction rests on one elementary change of viewpoint, and it is
worth presenting slowly, because everything after it is bookkeeping.

\subsection{A distribution is a rule assigning returns to ranks}

A probability law for the return $X$ is usually pictured as a density: how
much probability sits near each return level.  Invert the picture.  Order
all possible outcomes from worst to best and index them by their
\emph{rank} $u\in(0,1)$: $u=0.01$ is the first percentile, $u=0.5$ the
median, $u=0.99$ the ninety-ninth.  The law is then completely described by
the rule that assigns to each rank the return realized there,
\begin{equation}
u\;\longmapsto\;Q(u),
\label{eq:Qdef}
\end{equation}
the \emph{quantile function}, the inverse of the CDF.  If $U$ is uniform on
$(0,1)$ then $X=Q(U)$ has exactly the law described --- this is the
inverse-CDF construction used by every random-number generator.  Two facts
make the rank picture attractive for our problem.  First, the horizontal
axis is \emph{universal}: every law, however wild, is a rule on the same
interval $(0,1)$, and the uniform weight on ranks is the same for all of
them --- ranks are born normalized, so ``total probability one'' is not a
constraint but a tautology.  Second, the entire content of ``being a
probability law on returns'' collapses to a single property: $Q$ must be
\emph{increasing}.  Nothing else.

Monotonicity is a far better constraint to own than positivity of a second
derivative, because it is destroyed by one differentiation, not two.  Let
\begin{equation}
q(u)=Q'(u)
\label{eq:qdef}
\end{equation}
be the \emph{quantile density}: the spacing of returns per unit of rank.
Then $Q$ increasing $\iff$ $q>0$ $\iff$ $\log q$ is an arbitrary real
function.  Modeling $\log q$ therefore removes the last pointwise
constraint.  The quantity $\log q$ has a physical reading that will guide
every design choice below: it is the \emph{log-speed} at which the return
axis is traversed as the rank advances.  Where the speed is low, ranks are
packed into a narrow band of returns --- high probability density; where
the speed is high, a few ranks are stretched across a wide band ---
low density.  Indeed, differentiating $F_X(Q(u))=u$,
\begin{equation}
f_X\big(Q(u)\big)=\frac{1}{q(u)} :
\label{eq:reciprocal}
\end{equation}
density is the reciprocal of speed.  A model of $\log q$ is a model of the
distribution in which positivity of the density can never fail, for the
same reason $e^{t}$ can never be negative.

\subsection{The two singularities, and the coordinate that removes both}

The speed $q$ cannot be an arbitrary \emph{bounded} function, because an
unbounded law must escape to $\pm\infty$ as $u\to0$ or $1$: the integral of
$q$ over $(0,1)$ is the total return range, infinite for any law with
unbounded support.  So $\log q$ must diverge at both endpoints, and the
divergence is not a nuisance to be squashed but a structure to be
normalized.  Ask how fast the divergence is for the benchmark laws.  For a
standard Gaussian, $q(u)=1/\phin(\PhiN^{-1}(u))$, which blows up like
$1/(u\sqrt{2\log(1/u)})$ as $u\to0$ (using the Gaussian tail equivalent
$\PhiN(x)\sim\phin(x)/|x|$ as $x\to-\infty$, so
$\phin(\PhiN^{-1}(u))\sim u\,|\PhiN^{-1}(u)|\sim u\sqrt{2\log(1/u)}$) ---
essentially $1/u$, with a slowly varying correction.  For an exponential tail, $q(u)$ blows up like a
constant times $1/u$ exactly.  The pattern is general: the endpoint
divergences of $\log q$ are, at leading order, $-\log u$ at the left end
and $-\log(1-u)$ at the right, with the law's individual character riding
on a lower-order remainder.  For exponential tails the remainder is
\emph{bounded} --- it tends to the logarithm of the tail scale --- while
for the Gaussian it still drifts to $-\infty$, slowly.  Boundedness of the
remainder is precisely what will delimit the model's class
(\cref{prop:universal}).

This suggests separating skeleton from flesh,
\begin{equation}
\log q(u)\;=\;\underbrace{-\log u-\log(1-u)}_{\text{universal skeleton}}
\;+\;g(u),
\label{eq:skeleton}
\end{equation}
and modeling only the smooth part $g$.  There is a cleaner way to say the
same thing, and it is the form the whole computational apparatus will use.
Change the rank coordinate to \emph{log-odds},
\begin{equation}
z=\logit u=\log\frac{u}{1-u},
\qquad
u=\Lambda(z)=\frac{1}{1+e^{-z}},
\label{eq:logodds}
\end{equation}
so that $z=0$ is the median and $z=\pm\log 9\approx\pm2.20$ are the 90th
and 10th percentiles: equal steps in $z$ multiply the odds $u/(1-u)$ by a
constant factor.  The derivative of the map is
$\dd u/\dd z=\Lambda(z)\big(1-\Lambda(z)\big)=\rho(z)$, the logistic
density, which decays like $e^{-|z|}$ in both directions.  Now compute the
speed of the return assignment \emph{per unit of log-odds}:
\begin{equation}
\frac{\dd Q}{\dd z}
=q(u)\,\frac{\dd u}{\dd z}
=\frac{e^{g(u)}}{u(1-u)}\cdot u(1-u)
=e^{g(u)} .
\label{eq:speed}
\end{equation}
Both endpoint singularities cancel identically.  In the log-odds coordinate
the quantile of any law in our class is the integral of a smooth, strictly
positive, \emph{bounded} function --- the exponential of the very $g$ we
chose to model.  The coordinate $z$ does for ranks what the logarithm does
for prices: it moves the interesting boundary behavior to a place where it
is linear.

\subsection{The model in one sentence}

Assemble the pieces.  Take $Z$ a standard logistic variable (that is,
$Z=\logit U$ for $U$ uniform --- CDF $\Lambda$, density $\rho$), and define
the return by transporting $Z$ through the integrated speed:
\begin{equation}
X=x(Z),
\qquad
x(z)=m+\int_0^{z}e^{g(\Lambda(t))}\,\dd t ,
\label{eq:transport_heur}
\end{equation}
with $g$ a polynomial of modest degree in the rank and $m$ the scalar that
makes $\E[e^{X}]=1$.  The map $x$ is strictly increasing whatever the
polynomial, so $X$ always has a genuine law; the density is positive by
\eqref{eq:reciprocal}; the mass is normalized because ranks are; the mean
condition costs one scalar.  What the polynomial's two endpoint values
control --- the asymptotic speeds $e^{g(0)}$ and $e^{g(1)}$ --- turns out
to be exactly the exponential decay scales of the two return tails, and the
single validity condition of the entire construction will be the
requirement that the right tail decay fast enough to give $e^{X}$ a finite
mean.  \Cref{fig:transport} draws the three objects for the simplest
possible choice, $g$ constant: the logistic score with its percentile
marks, the resulting affine transport with the strike roots of two options
marked on it, and the flat-at-the-money smile it prices.  The next section
states the definition precisely and dispatches the constant-speed case in
closed form.

\begin{figure}[htbp]
\centering
\paperfig{\textwidth}{fig_transport.pdf}
\caption{The transport view for a constant-speed slice.  Panel A: the
logistic score density $\rho(z)$ on the log-odds line; the marked ticks are
the 1st, 10th, 50th, 90th, and 99th percentiles --- every percentile of an
unbounded law at a finite coordinate.  Panel B: the transport map $x(z)$,
an exact straight line when the log-speed $g$ is constant; the dots mark
the strike roots $z_k$ solving $x(z_k)=k$ for the at-the-money strike
$k=0$ and one out-of-the-money strike --- the two options priced by hand
on this very slice in \cref{sec:ledger}.  Panel C: the implied-volatility smile of this slice ---
flat at the money and gently convex, because exponential return tails carry
more far-strike value than the log-normal's, which the Black formula can
only express as higher implied volatility.}
\label{fig:transport}
\end{figure}
